% ============================================================================
% Sección 4: Implementación
% ============================================================================
\section{Implementación}

\subsection{Estructura del proyecto}

La librería \texttt{search-library} se organiza siguiendo las convenciones estándar de paquetes Python con el layout \texttt{src/}:

\begin{lstlisting}[language={}, caption={Estructura del proyecto}, label={lst:structure}]
search-library/
  src/search_library/
    algorithms/
      astar.py
    core/
      nodes.py
      problem.py
      result.py
      state.py
    graph/
      edges.py
      graph.py
    grid/
      grid.py
      grid_search.py
    heuristics/
      base.py
      euclidean.py
      manhattan.py
    exceptions/
      exceptions.py
    utils/
      helpers.py
\end{lstlisting}

\subsection{Módulo Core: Abstracciones fundamentales}

\subsubsection{Interfaz SearchProblem}

La clase abstracta \texttt{SearchProblem[T]} define el contrato que todo problema de búsqueda debe implementar, parametrizada por el tipo genérico $T$ del estado:

\begin{lstlisting}[caption={Interfaz abstracta del problema de búsqueda}, label={lst:problem}]
class SearchProblem(ABC, Generic[T]):
    @abstractmethod
    def initial_state(self) -> T: ...

    @abstractmethod
    def is_goal(self, state: T) -> bool: ...

    @abstractmethod
    def successors(self, state: T) 
        -> list[tuple[T, float]]: ...

    def heuristic(self, state: T) -> float:
        return 0.0  # Dijkstra por defecto
\end{lstlisting}

Este diseño permite que el algoritmo A* opere sobre cualquier dominio sin conocer la representación interna del estado.

\subsubsection{Nodo de búsqueda}

El \texttt{Node[T]} encapsula la información necesaria para el algoritmo, incluyendo la reconstrucción del camino:

\begin{lstlisting}[caption={Estructura del nodo de búsqueda}, label={lst:node}]
@dataclass(order=False)
class Node(Generic[T]):
    state: T
    parent: Node[T] | None = None
    g_cost: float = 0.0
    h_cost: float = 0.0
    f_cost: float = field(init=False)

    def __post_init__(self) -> None:
        self.f_cost = self.g_cost + self.h_cost

    def reconstruct_path(self) -> list[T]:
        path: list[T] = []
        current: Node[T] | None = self
        while current is not None:
            path.append(current.state)
            current = current.parent
        path.reverse()
        return path
\end{lstlisting}

\subsubsection{Resultado de búsqueda}

El resultado es un \texttt{dataclass} inmutable (\texttt{frozen=True}) que contiene la solución completa:

\begin{lstlisting}[caption={Contenedor de resultados}, label={lst:result}]
@dataclass(frozen=True)
class SearchResult(Generic[T]):
    path: list[T] = field(default_factory=list)
    total_cost: float = 0.0
    nodes_explored: int = 0
    success: bool = False
    explored_states: frozenset[T] = field(
        default_factory=frozenset
    )
\end{lstlisting}

\subsection{Algoritmo A*: Implementación central}

La implementación de A* emplea un heap binario (\texttt{heapq}) como cola de prioridad y un conjunto cerrado para evitar reexpansiones:

\begin{lstlisting}[caption={Implementación del algoritmo A*}, label={lst:astar}]
class AStarSearch(Generic[T]):
    def __init__(self, problem: SearchProblem[T]):
        self._problem = problem

    def search(self) -> SearchResult[T]:
        initial = self._problem.initial_state()
        h_cost = self._problem.heuristic(initial)
        start_node = Node(
            state=initial, g_cost=0.0, h_cost=h_cost
        )
        # (f_cost, counter, node) - counter rompe
        # empates con orden FIFO
        counter = 0
        open_list = []
        heapq.heappush(
            open_list, 
            (start_node.f_cost, counter, start_node)
        )
        g_scores = {initial: 0.0}
        closed_set: set[T] = set()

        while open_list:
            _, _, current = heapq.heappop(open_list)
            if current.state in closed_set:
                continue
            closed_set.add(current.state)
            if self._problem.is_goal(current.state):
                path = current.reconstruct_path()
                return SearchResult(
                    path=path,
                    total_cost=current.g_cost,
                    nodes_explored=len(closed_set),
                    success=True
                )
            for succ, cost in self._problem.successors(
                current.state
            ):
                if succ in closed_set:
                    continue
                tent_g = current.g_cost + cost
                if tent_g < g_scores.get(
                    succ, float("inf")
                ):
                    g_scores[succ] = tent_g
                    h = self._problem.heuristic(succ)
                    node = Node(
                        state=succ, parent=current,
                        g_cost=tent_g, h_cost=h
                    )
                    counter += 1
                    heapq.heappush(
                        open_list,
                        (node.f_cost, counter, node)
                    )
        return SearchResult(success=False)
\end{lstlisting}

\subsection{Módulo Graph}

El módulo de grafos proporciona una estructura de datos para grafos ponderados, tanto dirigidos como no dirigidos:

\begin{lstlisting}[caption={Grafo ponderado con adaptador a SearchProblem}, label={lst:graph}]
class Graph(Generic[T]):
    def __init__(self, directed: bool = True):
        self.directed = directed
        self._adjacency: dict[T, list[tuple[T, float]]] = {}

    def add_edge(self, source, target, weight=1.0):
        self._adjacency[source].append((target, weight))
        if not self.directed:
            self._adjacency[target].append(
                (source, weight)
            )

    def to_search_problem(self, start, goal, heuristic):
        return GraphSearchProblem(
            self, start, goal, heuristic
        )
\end{lstlisting}

\subsection{Módulo Grid}

El módulo de grids implementa un tablero 2D con soporte para obstáculos, costos variables y movimiento configurable:

\begin{lstlisting}[caption={Grid 2D con obstáculos y costos variables}, label={lst:grid}]
class Grid:
    def __init__(self, rows, cols, *,
                 allow_diagonal=False,
                 default_cost=1.0):
        self.rows = rows
        self.cols = cols
        self.allow_diagonal = allow_diagonal
        self._obstacles: set[Position] = set()
        self._costs: dict[Position, float] = {}

    @classmethod
    def from_matrix(cls, matrix, *,
                    obstacle_value=1,
                    allow_diagonal=False):
        """Crea grid desde matriz 2D."""
        ...

    def get_neighbors(self, position):
        """Retorna vecinos transitables y costos."""
        directions = (EIGHT_DIRECTIONS 
            if self.allow_diagonal 
            else FOUR_DIRECTIONS)
        ...
\end{lstlisting}

\subsection{Módulo Heuristics}

Las heurísticas siguen el patrón Strategy, permitiendo su intercambio transparente:

\begin{lstlisting}[caption={Heurísticas: Manhattan y Euclidiana}, label={lst:heuristics}]
class Heuristic(ABC, Generic[T]):
    @abstractmethod
    def estimate(self, state: T, goal: T) -> float:
        """Nunca sobreestima el costo real."""

class ManhattanHeuristic(Heuristic[tuple[int,int]]):
    def estimate(self, state, goal) -> float:
        return float(
            abs(state[0]-goal[0]) + 
            abs(state[1]-goal[1])
        )

class EuclideanHeuristic(Heuristic[tuple[int,int]]):
    def estimate(self, state, goal) -> float:
        dx = state[0] - goal[0]
        dy = state[1] - goal[1]
        return math.sqrt(dx*dx + dy*dy)
\end{lstlisting}

\subsection{Diagrama de flujo del algoritmo}

La Figura~\ref{fig:flowchart} ilustra el flujo de ejecución del algoritmo A* implementado.

\begin{figure}[H]
\centering
\begin{tikzpicture}[
    node distance=1cm,
    startstop/.style={rectangle, rounded corners, draw=black, fill=red!10, text width=2.5cm, text centered, minimum height=0.7cm, font=\scriptsize},
    process/.style={rectangle, draw=black, fill=blue!10, text width=3cm, text centered, minimum height=0.7cm, font=\scriptsize},
    decision/.style={diamond, draw=black, fill=green!10, text width=1.8cm, text centered, minimum height=0.5cm, aspect=2, font=\scriptsize},
    arrow/.style={-{Stealth[length=2mm]}, thick}
]
    \node[startstop] (start) {Inicio};
    \node[process, below=of start] (init) {Inicializar open con $s_0$};
    \node[decision, below=of init] (empty) {open vacío?};
    \node[startstop, right=2.5cm of empty] (fail) {Fallo};
    \node[process, below=of empty] (pop) {Extraer $n$ con menor $f(n)$};
    \node[decision, below=of pop] (goal) {$n \in G$?};
    \node[startstop, right=2.5cm of goal] (success) {Retornar camino};
    \node[process, below=of goal] (expand) {Expandir sucesores};
    \node[process, below=of expand] (update) {Actualizar $g$, $f$ en open};
    
    \draw[arrow] (start) -- (init);
    \draw[arrow] (init) -- (empty);
    \draw[arrow] (empty) -- node[above] {Sí} (fail);
    \draw[arrow] (empty) -- node[right] {No} (pop);
    \draw[arrow] (pop) -- (goal);
    \draw[arrow] (goal) -- node[above] {Sí} (success);
    \draw[arrow] (goal) -- node[right] {No} (expand);
    \draw[arrow] (expand) -- (update);
    \draw[arrow] (update.west) -- ++(-1.5,0) |- (empty.west);
\end{tikzpicture}
\caption{Diagrama de flujo del algoritmo A*}
\label{fig:flowchart}
\end{figure}
